\section{Introduction}\label{sec:introduction}
    Credit risk prediction involves determining whether an applicant will default on a future financial obligation \cite{zhu2025hybrid}, and is therefore a fundamental task in the financial sector. %In the literature, this topic is often referred to using different terms such as credit rating, credit scoring, loan evaluation, credit bond, and creditworthiness \cite{yeo2025comprehensive, lu2025llm, yu2025hybrid}. 
    Credit risk prediction is commonly considered as a supervised classification task; the goal is to assign each applicant to a class (e.g., default or non-default), rather than estimating a continuous predictive risk value. Credit risk prediction directly impacts the stability of banks and financial institutions, where accurate and explainable predictions are necessary to ensure regulatory compliance \cite{zhao2024improved, nguyen2025comparative}. Increasingly, financial institutions have implemented evaluations and management processes that employ pattern recognition and machine learning algorithms for credit risk prediction \cite{weber2024applications, yeo2025comprehensive}. As these predictors become more widely adopted, improving default prediction enables more accurate credit risk prediction and reduces losses \cite{lu2025llm, hartomo2025novel, yu2025hybrid}.\\

    Recent research in credit risk prediction has prioritized improving the performance of credit risk prediction algorithms based on machine learning and pattern recognition through the use of oversampling or undersampling methods, thereby reducing algorithmic bias towards the majority class \cite{adiputra2025effectiveness, ali2025explainable, kandi2025enhancing, darwish2025intelligent, hu2025combining}; the use of feature selection methods for identifying the most relevant features of a dataset to improve the performance of prediction algorithms \cite{zhu2025hybrid, gasmi2025features, zou2025variable, wang2024feature}; and the use of explainable AI techniques \cite{ali2025explainable, lin2025credit, zhang2025method, han2024nate} that allow users to comprehend and trust the predictions provided by credit risk prediction algorithms, which makes the predictions of the algorithm to comply with regulations.\\
        
    Machine learning algorithms have demonstrated good predictive performance, and decision trees are among the most effective for credit risk prediction \cite{weber2024applications, chang2024prediction, aruleba2024effective, zhang2025blockchain, xia2025confidence, wu2024credit, zou2025interpretable}. However, the need for explanations remains a significant limitation, particularly in the financial sector, as financial institutions and regulators require justification for their decisions \cite{sowmiya2024credit, ravi2025explainable}. Several authors have addressed the lack of explainability in credit risk prediction \cite{lin2025credit, zandi2025attention,sun2025enhancing}. \newline
    
    Explainability techniques, also known as XAI (eXplainable Artificial Intelligence), aim to make machine learning prediction processes transparent and understandable to humans. These techniques are important in areas such as credit risk prediction, where justifying and explaining the predictions is essential in enabling financial institutions to make more informed decisions. Several recent studies have applied SHAP \cite{lundberg2017unified} and LIME \cite{ribeiro2016should} to credit risk prediction to generate explanations for individual predictions \cite{garcia2025machine, baisholan2025fraudx, alonso2025should}. SHAP assigns to each feature a value that represents how much it contributed to the prediction. These values help explain the role each feature played in the prediction. LIME builds a simple interpretable predictor that approximates the behaviour of the original predictor by modifying a specific instance and observing the changes in the prediction. The result is an explanation that highlights the features contributing most to the prediction, showing how each affects the prediction in a positive or negative direction. However, these techniques focus on individual features and often overlook how several features jointly contribute to the prediction. \\

    In credit risk prediction, a predictor's output depends on multiple features rather than individual features. Therefore, explainable predictors should consider multiple features to explain the prediction \cite{zhang2025method, do2024explainable}. In response to this need, this paper proposes a credit risk prediction algorithm that generates explanations based on contrast patterns, which are conjunctions of conditions of the form $feature \# value$ that are frequent in a class but rarely appear in the opposite class. Since contrast patterns can be extracted from decision-tree-based, we propose using decision-tree-based predictors to predict credit risk, which have been reported  to perform competitively compared to the most successful approaches. In this way, the proposed algorithm leverages the affinity and synergy between decision-tree-based predictors, and contrast patterns extracted from decision treed. \newline
    
    Unlike other explainable algorithms that merely highlight the importance of individual features, our method explains predictions based on contrast patterns present in the training data and shared with the new applicant. Experiments conducted on both public and private financial datasets show that the proposed algorithm achieves competitive predictive accuracy while offering instance-level explanations based on the contrast patterns shared with other clients in the class, unlike the current state-of-the-art, which explains based on the single features contributing most to the prediction.\\
	
    %The remainder of this paper is organized as follows. Section \ref{sec:relatedworks} presents the related work. Section \ref{sec:explainablealgorithm} describes the proposed algorithm. Section \ref{sec:experimentsresults} presents the experimental results. Finally, Section \ref{suc:conclusion} presents our conclusions and future work directions.
    The remainder of this paper is organized as follows. Section 2 presents the related work. Section 3 describes the proposed algorithm. Section 4 presents our experimental results and a discussion about them. Finally, Section 5 presents our conclusions and future work directions.
    